% !TEX root = ../Projektdokumentation.tex
\section{Durchführung und Ergebnisse}
\label{sec:DurchfuehrungUndErgebnisse}
\subsection{Architektur und Zielplattform}
\label{sec:ArchitekturUndZielplattform}
Die Anwendung wurde als Webanwendung auf Basis von Python~3.12, FastAPI und
SQLModel geplant. Der Zugriff erfolgt über einen Browser, sodass auf den
Arbeitsplätzen keine eigene Client-Installation notwendig ist. Das passt zum
internen Einsatz, weil Administratoren die Anwendung zentral aufrufen können
und nicht mehr direkt auf der Shell arbeiten müssen.

Die Anwendung ist in vier Bereiche gegliedert: Datenmodell, Fachlogik, Web-API
und Templates. Im Quellcode entsprechen diese Bereiche den Ordnern
\Code{app/models}, \Code{app/services}, \Code{app/web/api} und
\Code{app/web/templates}. Dadurch ist klar, wo Daten gespeichert, Regeln
geprüft und Inhalte angezeigt werden. Diese Trennung erleichtert die Wartung
und die Tests, weil die Regeln nicht in einzelnen Masken verteilt sind.
Ein Komponentendiagramm ist im
\Anhang{app:Komponentendiagramm} dargestellt.

\subsection{Entwurf von Datenmodell und Rechtestruktur}
\label{sec:DatenmodellUndRechte}
Das Datenmodell umfasst Firmen, Benutzer,
Mail-Adresse und IP-Adresse. Ergänzt werden diese durch Tabellen für Audit- und
Login-Daten. Damit lässt sich später nachvollziehen, welche Änderungen
durchgeführt wurden und welche Anmeldevorgänge sicherheitsrelevant
waren.

Im Berechtigungskonzept wurden globale und firmenbezogene Rechte
getrennt. Root- und Admin-Rollen dürfen administrative Aufgaben für mehrere
Firmen ausführen. Einfache Benutzer bleiben dagegen auf die
Firmen beschränkt, denen sie zugeordnet sind. Diese Trennung ist
wichtig, damit Kundendaten nicht versehentlich firmenübergreifend
sichtbar oder bearbeitbar werden.
Das Datenbankmodell ist im \Anhang{app:Datenbankmodell} dokumentiert.

\subsection{Entwurf der Benutzeroberfläche}
\label{sec:EntwurfBenutzeroberflaeche}
Die Benutzeroberfläche wurde bewusst einfach gehalten. Im Mittelpunkt stand,
dass Daten schnell und mit möglichst wenigen Fehlern erfasst werden können.
Die Oberfläche wurde mit einem festen Navigationsbereich und einem zentralen
Inhaltsbereich entworfen. Bei einem Wechsel zwischen Dashboard, Firmen,
Benutzern, Mail-Adressen, IP-Adressen und Audit-Log soll nur der
Inhaltsbereich ausgetauscht werden, während Navigation und Grundlayout
erhalten bleiben. Dadurch muss nicht die gesamte Seite neu aufgebaut werden,
was die Bedienung bei Bereichswechseln beschleunigt und unnötige
Datenübertragungen reduziert. Zur Abstimmung des grundsätzlichen Aufbaus wurde vor der
Umsetzung ein Mockup erstellt. Dieses zeigt die geplante Navigation und die
tabellarische Pflegeansicht und ist im \Anhang{app:ProjektMockup} dargestellt.

Fehler werden direkt an den betroffenen Eingabefeldern angezeigt. Dadurch
fallen falsche E-Mail-Adressen, ungültige IP-Adressen oder unvollständige
Stammdaten sofort auf und müssen nicht später manuell gesucht werden. Die
einheitliche Bedienung soll den Einstieg erleichtern und fehlerhafte Eingaben
reduzieren.

\subsection{Sicherheitskonzept und Auditierung}
\label{sec:Sicherheitskonzept}
Neben der Bedienung wurde ein Sicherheitskonzept entworfen. Dazu
gehören Session-Login, Passwort-Hashing, Login-Audit, Begrenzung wiederholter
Fehlanmeldungen sowie kontobezogene Sperrmechanismen. Diese Funktionen sollen
verhindern, dass Verwaltungsfunktionen ohne Berechtigung genutzt werden.

Das Audit-Log wurde als eigener Teil der Anwendung vorgesehen. Jede relevante
Änderung an Firmen, Benutzern, Mail-Adressen und IP-Adressen wird mit
Zusatzinformationen protokolliert. Dadurch können fehlerhafte Pflegevorgänge
nachvollzogen und bewertet werden. Für ausgewählte Szenarien ist ein Rollback
vorgesehen, damit unbeabsichtigte Änderungen kontrolliert
zurückgenommen werden können.

\subsection{Schnittstellen und Datenfluss}
\label{sec:SchnittstellenUndDatenfluss}
Die Anwendung stellt sowohl HTML-basierte Oberflächen als auch API-Endpunkte
bereit. Dadurch können Daten direkt über den Browser gepflegt werden. Die
verwalteten Mail- und IP-Informationen stehen gleichzeitig strukturiert für
nachgelagerte Verarbeitungsschritte bereit. So müssen die Daten nicht doppelt
gepflegt werden und können einfacher weiterverarbeitet werden.

Für die Anbindung nachgelagerter Postfix-Prozesse wurde ein dateibasierter
Export vorgesehen. Pro Organisation soll unter \Code{postfix/<Organisationsname>/}
je eine Datei \Code{mail\_addresses.txt} und \Code{ip\_addresses.txt}
bereitgestellt werden. Exportiert werden dabei nur aktive Einträge. Damit kann
ein externer Folgeprozess die gepflegten Daten übernehmen, ohne selbst
auf die Anwendungsdatenbank zugreifen zu müssen.

Im Datenfluss beginnt jeder Pflegevorgang mit der Authentifizierung des
Benutzers. Danach prüft die Anwendung Rolle und Firmenzuordnung,
validiert die Eingaben, schreibt die Änderung in die Datenbank und erzeugt bei
relevanten Aktionen einen Audit-Eintrag. So wird bei jedem Vorgang
geprüft, ob der Benutzer die Änderung durchführen darf und ob sie später
nachvollzogen werden kann.
Bei Änderungen an Organisationen, Mail-Adressen oder IP-Adressen wird
anschließend der organisationsbezogene Postfix-Export aktualisiert, damit
Oberfläche, Datenbank und bereitgestellte Exportdateien konsistent bleiben.
